\textbf{Proof.}
Fix \(\eta\in\mathfrak C_D(\mathbf M)\).  Put \(E=\mathcal G_\eta[V]\), let
\[
H=\{\operatorname{pa}_{\mathcal G_\eta}(s):s\in S_\eta\},
\qquad
T_k=I_\eta^{(k)}\quad(k\in\mathcal K_+),
\]
where the braces in the definition of \(H\) discard multiplicity.  By \(\operatorname{ass{:}selection\text{-}parents}\), every element of \(H\) is a subset of \(V\).  There are only \(2^D\) subsets of \(V\), so \(|H|\le 2^D\).

Construct \(\bar\eta\) by retaining \(E\), the ordered targets, and one binary childless selection node \(s_h\) with parent set \(h\) for each \(h\in H\).  This preserves finiteness, acyclicity, causal sufficiency among vertices in \(V\), binary selection, childlessness, observed selection parents, the selected-stratum convention, pre-treatment timing, the empty observational target, and soft-intervention modularity.  It remains to verify the graphical projections.

The set of observed ancestors of selection is unchanged, because in either candidate it is
\[
A=\operatorname{an}_E\!\left(\bigcup_{h\in H}h\right).
\]
For each regime, the affected set \(F_k=\operatorname{de}_E(T_k)\) is also unchanged.  Multiplicity of a selection node with a given parent set affects neither the assertion that \(i,j\) share a selected child nor any of the sets \(A,F_k\).  The observational adjacency formula, every interventional path-induced adjacency, every indicator adjacency, and every endpoint-orientation case in \(\texttt{lem:dai-certificate-mag-characterization}\) therefore has the same truth value for \(\eta\) and \(\bar\eta\).  Hence
\[
\mathcal M_{\eta}^{(k)}=\mathcal M_{\bar\eta}^{(k)}
\qquad\text{for every }k\in\mathcal K.
\]
Since \(\eta\) matches the supplied representatives, so does \(\bar\eta\), proving \(\bar\eta\in\mathfrak C_D(\mathbf M)\).  The tuple \((E,H,(T_k))\) is the required normalized certificate. \(\square\)